\documentclass[preprint,11pt]{elsarticle}
\usepackage[T1]{fontenc}
\usepackage{float,longtable}
\usepackage[utf8]{inputenc}
\usepackage{lmodern,geometry,microtype,graphicx,booktabs,amsmath,amssymb,amsthm,tabularx,enumitem}

\usepackage[colorlinks=true,allcolors=blue,pdfauthor={Meshaal Obaid Mouawad}]{hyperref}
\geometry{margin=25mm}
\newtheorem{proposition}{Proposition}
\newcommand{\code}[1]{\texttt{#1}}
\journal{Science of Computer Programming}
\begin{document}
\begin{frontmatter}
\title{Version Comparisons Preserve Governance Alert Decisions Across Evidence and Policy Revisions}
\author[msu]{Meshaal Obaid Mouawad\corref{cor}}
\ead{mmouawad@ece.msstate.edu}
\cortext[cor]{Corresponding author}
\address[msu]{Department of Electrical and Computer Engineering, Mississippi State University, Box 9571, Mississippi State, MS 39762, USA}
\begin{abstract}
A reviewer may defer a warning about one revision of a software component. If a later revision produces a finding, the review episode can be active even when an obsolete deferral hides the warning. We ask which version comparisons allow a monitor to reconstruct the warning at each event and which allow it to report why that decision was made. We specify a policy that checks both the evidence revision and the policy version named by each observation and deferral. With packet presence, observation values, times and expiry retained, three joint comparisons preserve the displayed state, episode activity and severity. They record whether either current version changed and whether each supplied observation or deferral matches both current versions. Six separate comparisons also preserve the reason label. Thirty constructed pairs of event sequences expose the loss of a decision when one version field and its comparison are unavailable. The full monitor matched all 987 base sequences from three held-out rule families sharing 329 temporal templates. Accepting obsolete deferrals preserved every active interval and changed 30 complete decision sequences. Across 4,935 constructed sequences, the three joint comparisons reproduced all 26,430 event decisions. A decoder that assigned every version change or mismatch to the evidence revision changed 585 reason labels. A separately encoded monitor in the PyContract runtime-verification library reproduced all decisions and reasons. The two interfaces let a source-review system preserve the warning alone or also preserve the policy's explanation.
\end{abstract}
\begin{keyword}
governance \sep runtime monitoring \sep provenance \sep temporal requirements \sep software verification \sep evidence validity
\end{keyword}
\end{frontmatter}


\section{Introduction}\label{sec:introduction}
A software reviewer may defer a warning for a particular revision of a component. A later finding requires two decisions. Does the finding warrant an active review episode, a period in which the issue qualifies for review? If it does, does the earlier deferral still authorize withholding the warning? A monitor can answer the first question correctly and still hide a warning the reviewer should see. We identify which version comparisons preserve these decisions under a specified review policy, using proofs over event sequences and a controlled replay.

In a constructed audit case, positive findings begin at one second and qualify the current revision for a warning at three seconds. The monitor also receives a deferral that expires at four seconds and names a different evidence revision. A monitor checking only its expiry hides the warning at three seconds. A monitor checking the authorized revision shows it. Both record the same active episode. The difference is in what the reviewer sees at the qualifying event.

Two established capabilities explain how this situation can arise. Nweke and Yeng tie compliance evidence to releases and record the policy under which it was evaluated \cite{nwekeyeng2026compliance}. Prometheus Alertmanager retains active alerts and can withhold their notifications under a matching temporary silence \cite{prometheus2026alertmanager}. Henzinger et al. formalize how histories assigned to the same monitor class receive one output, and Jevitha et al. construct abstractions of program executions for selected properties \cite{henzinger2022abstract,jevitha2024abstract}. For revision-sensitive warnings, the unresolved interface question concerns the required output: which distinctions among the current versions and the versions named by a finding or deferral must reach a monitor that displays the warning, and which must also reach a monitor that explains it?

We investigate this question using five source checks from a system that links key performance indicator (KPI) documentation to source artifacts \cite{mouawad2026architecture}. Each check supplies a finding about one artifact. An \emph{event} records a review time, the current evidence and policy versions, and any finding and deferral supplied at that time. A \emph{trace} is an ordered sequence of such events for one artifact and rule family. A deferral is the authorization to withhold a warning; \emph{suppression} is its effect on the displayed decision. The monitor uses successive events to open, sustain, clear or escalate an episode and to determine the warning shown at each event.

The evaluation follows three questions. \textbf{Q1:} Which version distinctions preserve exact decision sequences and their reasons? \textbf{Q2:} How does omitting evidence-version, deferral-version or evidence-age checks affect the displayed warning and the underlying episode? \textbf{Q3:} Do irrelevant metadata edits or equivalent changes of time unit alter the output? We connect these questions through three contributions:
\begin{itemize}[leftmargin=*,itemsep=3pt,topsep=3pt]
\item An executable applicability contract separates the evidence supporting a review episode from the authorization to suppress its warning.
\item Sufficiency proofs identify three joint comparisons for displayed state, activity and severity, and six separate comparisons that also preserve the specified reason. Constructed pairs expose the decision loss caused by deleting a version field without retaining its comparison.
\item An evaluation of complete decisions and active episodes exposes an obsolete-deferral error: accepting such deferrals preserves all 987 held-out base episode sequences and changes 30 complete decision sequences.
\end{itemize}
The paper develops the source-review context and closest work, specifies the contract, proves its information requirements, and evaluates the resulting decisions and explanations.

\section{Background and related work}
\subsection{Source findings and review episodes}
The source-documentation system is LEAP. It generates a finding from one source artifact at a time. The dispatcher labels each finding \code{REVIEW\_REQUIRED}. Five lexical rules supply the Boolean observations used here (Table~\ref{tab:recovery}). The episode timers and the binding of a deferral to a revision extend this snapshot interface. The companion source inventory identifies the adapters and their archived versions.

\begin{table}[ht]
\centering\small
\caption{Recovered source-review predicates used as Boolean observation adapters. The evaluated predicate is the stated lexical check.}\label{tab:recovery}
\begin{tabularx}{\linewidth}{@{}lX@{}}\toprule
Family & Source-evidence pattern \\\midrule
Privacy & Personal-identifier tokens without masking/encryption-related substrings.\\
Restricted execution & Tokens such as \code{eval}, \code{exec}, or \code{os.system}.\\
Audit completeness & Modification tokens with one or more missing audit-signature elements.\\
Data transfer & Personal-data and movement tokens without consent/agreement tokens.\\
Lineage signing & Lineage tokens without signature-related tokens.\\\bottomrule
\end{tabularx}
\end{table}

The five predicates scan text. For example, a masking-related comment can change the privacy predicate's output, and an \code{eval} token inside a comment can trigger the execution predicate. The new experiment fixes the finding interface and evaluates temporal decisions from its outputs. The reconstructed source adapter scans controlled snippets; it never executes those snippets. This preserves a precise separation between the inherited evidence detector and the newly developed lifecycle.

\subsection{Input information and preserved outputs}\label{sec:related-work}
An input representation preserves a monitor's output when histories that become indistinguishable require the same result. Henzinger et al. formalize this relationship for quantitative specifications by grouping observation prefixes into classes that remain equivalent under a common continuation. Each class receives one real-valued output. Their Proposition~6 bounds the difference between the two specified values by twice the permitted prompt error; exact agreement requires the values to coincide. They use these classes to relate output precision to the distinctions a monitor can retain \cite{henzinger2022abstract}. This provides the foundation for testing whether an interface merges histories that require different warnings.

Leucker et al. study a complementary representation of uncertainty: an abstract timed stream denotes the complete streams consistent with its gaps or imprecise values. Their abstract TeSSLa operators retain the precision obtainable by applying the corresponding concrete operator to those possibilities and then abstracting the results. For \code{lift}, this property requires a perfect abstraction of the underlying data function. Their analysis also shows why composing individually perfect operators can lose temporal information, and gives a time-aware construction that retains the relevant bounds \cite{leucker2019partial}. The representation and the operations applied to it jointly determine what a monitor can recover.

Jevitha et al. choose abstractions according to the property being checked. They derive finite-state models from multithreaded Java executions and reduce them using user-defined data and control abstractions. Predicate checks concern the concrete states represented by an abstract state; selected path conditions concern the connections retained between states. Their evaluation compares state counts and model-construction and verification times \cite{jevitha2024abstract}. The common question across these approaches is which distinctions the chosen output requires. In the present policy, an obsolete observation context and an obsolete observation policy can both produce an unresolved warning. A single mismatch indicator captures that decision, whereas the prescribed reason distinguishes its source.

Antignac et al. formalize a data preprocessor that preserves a program's output and give a procedure for constructing such preprocessors. Their ordering compares the input distinctions retained for a fixed output function \cite{antignac2017minimisation}. Here, the two output functions are the event decision and that decision with its prescribed reason; Propositions~3 and~4 establish concrete comparison encodings for the review lifecycle.

\subsection{Temporal monitoring and explanations}
Runtime-monitoring languages supply ways to express the transitions that consume those inputs. Convent et al. define transformations of timestamped streams, including delay-generated events, in TeSSLa \cite{convent2018tessla}. Baumeister et al. describe specifications over streams, derived streams and Boolean notification triggers in RTLola, with static analyses and a configurable interface to the observed system \cite{baumeister2026rtlola}. These mechanisms support the execution of a review policy after its observation and output contracts have been defined.

The event source can also determine the information available for verification. Chalupa et al.'s Vamos middleware preprocesses heterogeneous streams and represents observations lost under load through explicit gaps or summaries, such as event counts and timestamp ranges \cite{chalupa2025vamos}. Ferrando and Cardoso's partial-monitoring construction instead identifies when checking can stop or observations can be skipped because they cannot contribute to a conclusive verdict \cite{ferrando2025partial}. Both address which observations the monitor needs to process. A version-applicability comparison addresses a recorded observation whose content is available and whose authority refers to a particular evidence or policy revision.

Timing and changes to requirements add two further choices. Leucker et al. distinguish system time from event-arrival time when computing verdicts from asynchronous streams \cite{leucker2020streams}; the present evaluation specifies decisions at supplied event times. Zhou et al. analyze on/off-delay timers for stochastic anomaly detectors \cite{zhou2025delay}. Such timers determine when a condition persists long enough to change a decision. Carwehl et al. adapt runtime verification to changing requirements using timed-automata patterns and transfer relevant monitor state \cite{carwehl2023changing}. The review contract below specifies which evidence and policy changes instead reset that state.

Explaining a verdict introduces an additional output requirement. Zhang et al. use causation semantics to describe the relevance of incoming observations to an online Signal Temporal Logic verdict \cite{zhang2026caumon}. Lima et al. construct proof-tree explanations for metric first-order temporal logic and provide a verified checker; their implementation is WhyMon \cite{lima2024explainable}. The present monitor reports the contract condition selected at an event, such as an observation-policy mismatch. Its information question is whether a reduced interface retains the distinction needed to select that reason as well as the displayed decision.

\subsection{Versioned evidence and notification suppression}
The W3C PROV data model, edited by Moreau and Missier, represents revision and invalidation relationships \cite{moreau2013prov}. Operational assurance systems use version information to decide which evidence supports a finding or release. Muhammad et al. combine versioned policies and schema-checked evidence bundles in a deterministic pass/warn/block gate, evaluated on 1,000 synthetic scenarios and five benchmark model configurations \cite{muhammad2026audit}. Nweke and Yeng bind runtime evidence to controls, releases, environments and timestamps, with policy or configuration digests where required. Their four-release evaluation checks provenance and evidence completeness in two identity workflows \cite{nwekeyeng2026compliance}. Open Policy Agent, used in that implementation, records the policy-bundle revision with each decision identifier \cite{opa2026logs}.

Versioned records also support temporal checking. Li et al. record model versions with threshold triggers, observation windows, metrics and proposed mitigations in a clinical audit ledger \cite{li2026sentinel}. Ra\c{t}iu et al. check temporal constraints incrementally as related engineering artifacts change, using Object Constraint Language with linear temporal operators to report process deviations over artifact histories \cite{ratiu2026runtime}. These systems connect observations or findings to the objects and versions that give them meaning.

Suppressing a notification introduces a separate authorization. Prometheus Alertmanager stores active alerts and can mute notifications when their labels match a temporary silence \cite{prometheus2026alertmanager}. The distinction between stored activity and notification gives the review policy two relationships to check: whether the observation supports the current episode and whether the deferral applies to its warning. The constructed audit case uses that separation to expose a display error with an unchanged active interval.

Version-sensitive exceptions also occur in analysis tools. The \code{skill-audit} implementation binds a suppression to a fingerprint of the source, evidence, rule and scanner version; a changed fingerprint reactivates the finding. Its reports retain the finding and record the suppression separately \cite{skillaudit2026suppression}. These mechanisms establish revision-scoped exception handling. The present analysis determines which comparisons preserve the complete temporal decision and its specified reason.

Companion LEAP studies examine checks on generated KPI documentation and comparisons between source-pinned KPI definitions and their implementations \cite{mouawadjones2026checks,mouawad2026discrepancies}. Those checks supply findings about fixed source material. Here, successive revisions become inputs to a temporal policy and can change whether a finding or deferral applies. Table~\ref{tab:related-matrix} places this interface question alongside the abstractions, temporal mechanisms and versioned evidence objects established by the closest work. The following sections define the contract and establish sufficient comparisons for each output requirement.

{\small
\setlength{\tabcolsep}{4pt}
\renewcommand{\arraystretch}{1.12}
\setlength{\LTleft}{0pt}
\setlength{\LTright}{0pt}
\begin{longtable}{@{}>{\raggedright\arraybackslash}p{.235\linewidth}>{\raggedright\arraybackslash}p{.38\linewidth}>{\raggedright\arraybackslash}p{\dimexpr.385\linewidth-4\tabcolsep\relax}@{}}
\caption{Comparison of input abstractions, monitored objects, and output contracts established in the cited sources.}\label{tab:related-matrix}\\
\toprule
Approach & Input and established operation & Decision or evidential object \\
\midrule
\endfirsthead
\multicolumn{3}{l}{Table~\thetable\ (continued)}\\
\toprule
Approach & Input and established operation & Decision or evidential object \\
\midrule
\endhead
\midrule
\multicolumn{3}{r}{Continued on next page}\\
\endfoot
\bottomrule
\endlastfoot
RTLola \cite{baumeister2026rtlola} & Stream specifications and asynchronous inputs; derived streams, static analyses and executable monitor construction. & Stream values and Boolean notification triggers. \\
Vamos \cite{chalupa2025vamos} & Heterogeneous event sources; preprocessing, buffering and explicit gaps or summaries when observations are lost. & Events and loss information supplied to downstream monitors. \\
Abstract TeSSLa \cite{leucker2019partial} & Timed streams with gaps or imprecise values; abstract operators and analysis of compositional precision. & Sound abstract outputs; operator-level perfection and selected perfect compositions. \\
Abstract monitors \cite{henzinger2022abstract} & Equivalence classes of observation prefixes; precision and resource requirements for quantitative monitors. & Common outputs for indistinguishable histories; bounds relating prompt error to specification-value separation. \\
Property-preserving Java models \cite{jevitha2024abstract} & User-defined data and control abstractions of execution-state models. & Preserved predicates and path conditions for runtime verification on reduced models. \\
Data minimisation \cite{antignac2017minimisation} & A program's input domains and output function; construction and comparison of output-preserving preprocessors. & The program's output with reduced input distinctions. \\
CauMon \cite{zhang2026caumon} & Signal traces and temporal requirements; online monitoring with causation semantics. & A verdict identifying the relevance of incoming observations to satisfaction or violation. \\
WhyMon \cite{lima2024explainable} & Metric first-order temporal requirements and event traces; proof-tree construction with a verified checker. & Explanations certifying satisfaction or violation. \\
Changing-requirement monitoring \cite{carwehl2023changing} & Requirement changes in self-adaptive systems; adaptation through timed-automata patterns. & Updated verification behavior with transfer of relevant monitor state. \\
Audit-as-Code \cite{muhammad2026audit} & Versioned policy and schema-checked evidence bundles; deterministic assurance gates. & Pass/warn/block decisions with policy-linked evidence. \\
Compliance-as-Code \cite{nwekeyeng2026compliance} & Control, release, environment and time-bound evidence; provenance and fixed-window completeness checks. & Evidence-linked findings and machine-readable assessment artifacts. \\
Temporal engineering compliance \cite{ratiu2026runtime} & Changes to related engineering artifacts; incremental checking of temporal constraints. & Feedback on process deviations as artifact histories evolve. \\
Alertmanager \cite{prometheus2026alertmanager} & Active alerts, label matchers and temporary silences. & Notification suppression separated from stored active alerts. \\
\code{skill-audit} \cite{skillaudit2026suppression} & Source, evidence, rule and scanner fingerprints; exact suppression matching with expiry. & Findings retained in reports with separate suppression annotations. \\
Present study & Current, observation and suppression bindings; interface projections and a specified review lifecycle. & Necessary applicability distinctions, sufficient comparison representations, complete decision sequences and active-episode intervals. \\
\end{longtable}
}

\section{The review contract}
The recovered rules supply findings about individual source records. The lifecycle below uses those findings, the current revisions, and elapsed time to decide whether a review episode is active and what to display at each event.

\subsection{Policy and event contract}
One trace concerns one source artifact and one rule family. Each event contains a time, a unit, a current relevant-evidence context $c$, and a current policy version $p$. An optional observation contains a Boolean finding or a conflict marker, the context and policy under which it was produced, and its observation timestamp. An optional suppression contains its authorized context and policy and an expiry time. Context refers to the relevant evidence revision; display metadata lies outside that identity. The policy supplies qualification duration $D$, clearing duration $C$, escalation duration $E$, and maximum evidence age $F$. In the audit example, the current and observation versions are $(c_1,p_1)$; the deferral names $(c_0,p_1)$ and expires at four seconds.

The contract distinguishes two decisions that a source-review system must make independently. An observation can justify opening an episode, while a deferral authorizes withholding its notification. The authorities for those two decisions can name different revisions. Release-linked evidence contracts, revision-sensitive suppressions and label-scoped silences provide independent precedents for these choices \cite{nwekeyeng2026compliance,skillaudit2026suppression,prometheus2026alertmanager}. We combine them to test the consequence of accepting a deferral that names the wrong revision. The dwell and expiry values make boundary cases executable; the three- and six-comparison proofs use the specified comparisons and transition order for any admissible duration parameters.

The experimental configuration is $D=C=F=2$ seconds and $E=4$ seconds. Qualification, clearing, escalation, and freshness boundaries are inclusive. Suppression applies strictly before its expiry. Seconds and milliseconds are normalized before comparison. Required timestamps and durations must be finite, durations nonnegative, and event times strictly increasing. Malformed input raises an execution failure. These choices form a new, explicitly specified review policy.

An observation is valid when its context and policy match the current event, its age lies in $[0,F]$, and its value is Boolean. Missing observations, conflicting values, future timestamps, stale observations, and version mismatches are unresolved states. Validation has a fixed precedence: missing, context mismatch, policy mismatch, future, stale, then conflict. This makes a multiply invalid packet's reported reason deterministic.

Figure~\ref{fig:method-pipeline} shows how a source-review finding becomes a lifecycle decision. The interface combines that finding with the current evidence and policy identities, time, and any suppression. The monitor first applies version-change resets, then checks the observation and advances episode state, and finally determines the displayed action. This order preserves the distinction between the existence of an active episode and the warning presented at a particular event.

\begin{figure}[H]
\centering
\includegraphics[width=0.88\linewidth]{method_pipeline.pdf}
\caption{Event processing under the version-bound review policy. After schema and time checks, a version change resets activity and timers. Valid observations update qualification and clearing before the monitor determines escalation and applies suppression. Unresolved observations display \code{UNKNOWN} and preserve any episode still active after the reset check. Suppression requires an active episode, matching context and policy, and event time $t$ strictly before expiry. At exact expiry the underlying episode again determines the valid display. Figure preparation and verification code used OpenAI ChatGPT/Codex assistance; the original vector export records Matplotlib 3.10.8.}
\label{fig:method-pipeline}
\end{figure}

\subsection{Transitions and displayed decisions}
The incremental monitor stores the previous context/policy pair, qualification and clearing start times, activation time, and an active flag. At a context or policy change it first resets the active episode and all timers. A valid positive finding starts or continues qualification; the first event at which its duration reaches $D$ activates the episode. A valid negative finding starts or continues clearing; the first event at which its duration reaches $C$ ends the episode. A valid negative finding resets qualification, and a valid positive finding resets clearing. Unresolved observations discard qualification and clearing runs while retaining an already active episode and its activation time.

For valid evidence, an active episode becomes escalated at the first event at least $E$ seconds after activation. If the observation is valid and the episode is active, an unexpired suppression that matches both current versions changes the displayed state to \code{SUPPRESSED}. The episode remains active and its escalation clock continues. An unresolved observation displays \code{UNKNOWN}, including when an active episode is retained. Other states are \code{CLEAR}, \code{PENDING}, \code{ACTIVE}, and \code{ESCALATED}. Escalation has severity HIGH; active decisions use the recovered rule severity; suppressed and inactive states have severity NONE; unresolved states have severity UNKNOWN. Thus a displayed state, an active flag, and a severity jointly define an event decision.

This distinction is important for evaluation. The active flag records the existence of an underlying review episode; the displayed state records the action presented at the current event. A suppression error can change the latter without changing the former. The implementation also reports a reason: the evidence-validity condition, a policy/context reset, suppression, or an ordinary state decision. This is a directly checked policy explanation.

\paragraph{Concurrent conditions and the reported reason}
The transition order and the reason-selection order are specified separately. Each event names one current context and one current policy. A change in either first resets episode activity and all timers. If both change, the reset label is \code{POLICY\_CHANGE}; a context-only change gives \code{CONTEXT\_CHANGE}. Observation validation then selects the first applicable condition in the ordered list \code{MISSING}, \code{CONTEXT\_MISMATCH}, \code{POLICY\_MISMATCH}, \code{FUTURE}, \code{STALE}, \code{CONFLICT}. A malformed packet fails schema validation before these decisions are produced.

At event index $i$, let $q_i$ be the selected observation-validity label, or empty for valid evidence, and let $r_i$ be the selected reset label, or empty when no reset occurs. Let $d_i$ denote the displayed state after the state update and suppression overlay. The reported reason is
\begin{equation}
\operatorname{reason}_i=
\begin{cases}
q_i, & q_i\ne\varnothing,\\
r_i, & q_i=\varnothing\ \land\ r_i\ne\varnothing,\\
\text{SUPPRESSION}, & q_i=r_i=\varnothing\ \land\ d_i=\text{SUPPRESSED},\\
\text{STATE}, & \text{otherwise},
\end{cases}
\end{equation}
The selected validity label takes diagnostic precedence even though a simultaneous version reset has already changed the underlying episode state. For example, a current-policy change with an observation still bound to the previous policy produces \code{UNKNOWN}, inactive activity and reason \code{POLICY\_MISMATCH}. If both observation bindings mismatch, the reason is \code{CONTEXT\_MISMATCH}. If matching evidence is both stale and conflicting, the reason is \code{STALE}.

For valid evidence, timer updates determine activity before escalation and suppression are displayed. Clearing at its boundary yields \code{CLEAR} even when escalation would otherwise become due at that event. Valid suppression overlays an active or escalated display while retaining the activation clock. At exact suppression expiry, the overlay ends. Unresolved evidence displays \code{UNKNOWN} with its validity reason, including when a valid suppression is present. When a valid reset and suppression coincide with immediate activation under zero qualification duration, the displayed state is \code{SUPPRESSED} and the reason is the reset label.

The evaluation includes a separately frozen set of concurrent-condition cases to check this priority across four implementations.

\section{Information requirements}
At an event with valid evidence and an active episode, two unexpired deferrals can require different displays if only one names the current revision. Erasing the deferral's version without retaining its comparison makes those cases look the same to the monitor. This simple pair motivates two questions: which information losses make the decision ambiguous, and which comparisons can replace raw version names while preserving the required outputs?

Proposition~1 applies the standard indistinguishability principle to histories grouped by a fixed input projection. Henzinger et al. give its quantitative counterpart for abstract monitors in Proposition~6 \cite{henzinger2022abstract}. Proposition~2 is the empirical form of the maximum-posterior decision rule under zero--one loss. For a trace drawn uniformly from the benchmark multiset, the normalized bound is the maximum probability of recovering the required sequence after observing the whole projected trace \cite[Definition~3]{smith2009foundations}. Propositions~3 and~4 then establish sufficient comparisons for the specified policy and separate preservation of a decision from preservation of its reason.

Let $X$ be the set of admissible complete traces, $Y$ the set of decision sequences, and $Z$ the set of projected traces. The function $G:X\rightarrow Y$ returns the reference sequence of displayed states, activity flags, and severity labels. The projection $P:X\rightarrow Z$ removes a specified field. A deterministic monitor restricted to this interface computes $M:Z\rightarrow Y$. Policy parameters and all remaining side inputs are included in $Z$.

\begin{proposition}[Indistinguishable evidence]
If $P(x_1)=P(x_2)$ and $G(x_1)\ne G(x_2)$, no deterministic function of the projected input is correct on both traces.
\end{proposition}
\begin{proof}
The two equal projected inputs produce one output under $M$. That output cannot equal two different required decision sequences.
\end{proof}

\begin{proposition}[Finite-benchmark bound]
For a finite multiset of $N$ traces, let $n(z,y)$ count traces with projected input $z$ and reference sequence $y$. Every deterministic function of that projection has at most
\begin{equation}
B(P)=\sum_z\max_y n(z,y)
\end{equation}
correct complete traces.
\end{proposition}
\begin{proof}
Within each projection class, the function returns one sequence and can match at most the largest reference-sequence group. Summation over disjoint projection classes gives the bound. A full-trace lookup function can attain it by choosing a most frequent label in each class. The same bound applies to causal monitors.
\end{proof}

Figure~\ref{fig:witness} illustrates a constructed observation-context witness. The two traces differ only in the observation's context after the first event. Current context, current policy, Boolean values, and times are equal. Removing the observation-context field makes the remaining inputs identical; one reference sequence activates and the other remains unresolved. A corresponding suppression witness distinguishes \code{ACTIVE} from \code{SUPPRESSED}, with activity retained in both cases.

\begin{figure}[H]
\centering\includegraphics[width=\linewidth]{necessity_witness.pdf}
\caption{A constructed observation-context witness. Boxes show the required displayed state at successive supplied events, equally spaced by event order. Both observations match the current context at 0\,s. From 2\,s onward only the upper row matches, so qualification produces \code{ACTIVE} there while the lower row displays \code{UNKNOWN}. Removing the observation-context field makes the remaining inputs identical. Figure preparation and verification code used OpenAI ChatGPT/Codex assistance; the original vector export records Matplotlib 3.10.8.}\label{fig:witness}
\end{figure}

Six field projections are examined separately: current context, current policy, observation context, observation policy, suppression context, and suppression policy. Each projection has one pair per rule family. Trace identifiers and intervention labels are excluded from the monitor interface. The analytical witness assay fixes base severity to High; six structural pairs are instantiated under five rule identifiers. The performance benchmark uses each recovered severity string. These inherited strings are categorical labels; HIGH escalation is a study-defined label, with no empirically ranked severity claim. The witness assay retains full inputs, canonical projected inputs, their hashes, and both reference sequences. It evaluates necessary information for this policy independently of the choice of implementation algorithm.

\begin{proposition}[Sufficiency of version comparisons]
For an admissible trace, the raw version names can be replaced by six Boolean comparisons per event: whether the current context and policy each changed since the preceding event, whether an observation's context and policy each match the current event, and whether a suppression's context and policy each match. The observation and suppression comparisons are defined when their respective packets are present. Given the remaining inputs, these comparisons suffice for the displayed decision, activity, severity, and reason at every event.
\end{proposition}
\begin{proof}
Initialize both change indicators to false. Suppose two encoded prefixes agree through event $i$ and have identical timers and activity at its start. The two change indicators determine the reset and its reason. The observation comparisons, time, and value determine validity and hence the same timer update. The suppression comparisons and expiry give the same displayed state. Activity, severity, and reason then agree, as do the timers for the next event. The common initial state establishes the induction.
\end{proof}
An evidence service can send these comparison results in place of version names. This preserves reason labels that distinguish evidence changes or mismatches from policy changes or mismatches. The next result asks which comparisons suffice when the receiving monitor reports only the decision components. 

\paragraph{An interface for decisions alone}
The following encoding uses $r_i$ as a Boolean change indicator. Let $k_i=(c_i,p_i)$ denote the current evidence-context and policy pair at event $i$. Let $r_i$ indicate that $k_i$ differs from $k_{i-1}$, with $r_1=0$. Let $v_i$ and $m_i$ indicate respectively that the observation and suppression, when present, name $k_i$. Presence, observation value, timestamps and expiry remain inputs. At most three Boolean comparisons therefore accompany each event.
\begin{proposition}
The triple $(r_i,v_i,m_i)$ suffices for the sequence of displayed state, activity, and severity under the stated policy. It need not determine the reported reason.
\end{proposition}
\begin{proof}
The monitor resets when $r_i$ is true. For a present observation, either kind of version mismatch clears the qualification and clearing runs and produces the same unresolved display; $v_i$ distinguishes that case from a matching observation. Given valid evidence and an active episode, a present deferral overrides the display exactly when $m_i$ holds and the expiry time is later than the event time. All other transitions use the retained time and value fields. Induction from the common initial state gives identical activity, display, and severity at each step. For the reason field, consider two otherwise identical first events with positive observations: one names an obsolete context under the current policy and the other names the current context under an obsolete policy. Both give $v_1=0$ and the same unresolved decision. The specified reasons are context mismatch and policy mismatch, respectively.
\end{proof}

\section{Controlled evaluation}
The paired-history assay and comparison re-encodings address Q1. The primary benchmark compares complete and reduced implementations to address Q2, and the metadata and time-unit transformations address Q3. Separate reference, schema, concurrent-condition and implementation checks verify the execution used to answer those questions.

\subsection{Benchmark and split}
The benchmark consists of newly constructed revision sequences over five recovered predicate families. Positive and negative snippets are checked by the recovered source predicates before being converted into observation packets. The original provenance links each rule family to its adapter and positive and negative snippets. Boolean observation packets in the historical archive retain their source snippet and hash. The offline replay starts from the resulting temporal packets. The historical fixtures check the recovered source rules. Each event trace used to evaluate the new lifecycle was constructed for this study.

For each rule family, the frozen generator defines 329 base designs: $4^4=256$ length-four sequences over positive, negative, conflicting and missing observations; three length-twelve controls; $11\times5=55$ intervention designs; and $5\times3=15$ temporal-boundary designs. The length-four sequences have event times $0,2,4,6$\,s. The controls comprise twelve negative observations, twelve positive observations, and an activation--clearing sequence whose values at $t=0,\ldots,11$\,s are one negative, seven positive and four negative observations.

Every intervention starts from that activation--clearing control. Each of the eleven operations in Table~\ref{tab:intervention-taxonomy} is applied at each index $j\in\{2,3,5,7,9\}$, yielding exactly 55 designs per family. Here the zero-based index equals event time in seconds. A single intervention means one selected operation: the four evidence-quality operations change one event, whereas version changes and suppression assignments affect the indicated suffix. The supplementary grid lists all 55 type--position pairs and their frozen trace identifiers.

\begin{table}[H]
\centering\small
\caption{Complete intervention taxonomy. Each row contributes five base designs per rule family, one at each intervention index. The ordinary context and policy are $c_1,p_1$; $c_0,p_0$ denote nonmatching identities and $c_2,p_2$ revised identities.}
\label{tab:intervention-taxonomy}
\begin{tabularx}{\linewidth}{@{}>{\raggedright\arraybackslash}p{0.26\linewidth}>{\raggedright\arraybackslash}X>{\raggedright\arraybackslash}p{0.12\linewidth}@{}}
\toprule
Operation & Exact modification & Extent \\
\midrule
Missing observation & Remove the observation packet at index $j$. & Event $j$ \\
Conflicting observation & Set its finding value to null. & Event $j$ \\
Stale observation & Set observation time to $j-3$\,s, giving age $3>F$. & Event $j$ \\
Future observation & Set observation time to $j+1$\,s. & Event $j$ \\
Revised context, old evidence & Set current context to $c_2$ and retain observation context $c_1$. & $i\ge j$ \\
Revised policy, old evidence & Set current policy to $p_2$ and retain observation policy $p_1$. & $i\ge j$ \\
Revised context, fresh evidence & Set both current and observation context to $c_2$. & $i\ge j$ \\
Revised policy, fresh evidence & Set both current and observation policy to $p_2$. & $i\ge j$ \\
Matching suppression & Attach suppression for $(c_1,p_1)$, expiring at $j+2$\,s. & $i\ge j$ \\
Old-context suppression & Attach suppression for $(c_0,p_1)$, expiring at $j+2$\,s. & $i\ge j$ \\
Old-policy suppression & Attach suppression for $(c_1,p_0)$, expiring at $j+2$\,s. & $i\ge j$ \\
\bottomrule
\end{tabularx}
\end{table}

The fifteen boundary designs test qualification, clearing, escalation, freshness and suppression at offsets $-0.5$, $0$ and $+0.5$\,s from the relevant boundary. They contain respectively two, four, three, two and three events per design. Suppression packets remain present after expiry in their intervention suffixes, allowing the policy to evaluate expiry explicitly.

Each base design has three representations: its original form, a variant with irrelevant display metadata added to every event, and an equivalent millisecond representation scaling event times, observation times and suppression expiry by 1,000. Across five families, the corpus therefore contains 1,645 design groups and $329\times3\times5=4,935$ traces. The four design categories contribute, respectively, 3,840, 45, 825 and 225 traces, comprising 15,360, 540, 9,900 and 630 events; the total is 26,430 events. The taxonomy audit compares the policy-input fields of every intervention representation with the frozen records after time-unit normalization.


The partition is by source-rule family. Privacy and restricted execution supply the 658 development base traces. Audit, transfer, and lineage supply 987 held-out base traces and 5,286 events, with 2,961 traces and 15,858 events after the two transformations. Simulated source identifiers and rule families are disjoint across the partition, and the three representations of each family-design group remain together. All five families instantiate the same 329 temporal templates. The source predicates supply Boolean inputs upstream, and the temporal evaluation checks contract conformance across these family instances. The sample size follows finite mechanism coverage.

\subsection{Reference and verification protocol}
A separate AI review agent implemented the reference from the written specification without viewing the incremental monitor. The reference reconstructs decisions from prefixes and contiguous valid observation runs. Before benchmark scoring, 29 hand-derived cases specified transitions, exact boundaries, invalid evidence, suppression and malformed inputs. The specification, implementations, analysis script, inputs and reference outputs were frozen before the held-out comparison; the manifest records their identities and freeze time.

Post-evaluation verification uses a separately versioned validation reference that requires an explicit time-unit field during schema validation. The archive retains this reference and the original one. Their outputs are compared on the complete benchmark and witness set. Twenty-eight malformed-input probes check the schema, while 24 hand-calculated concurrent-condition traces check the precedence rules in the review contract. The concurrent-case inputs and expectations were frozen before execution and scored in the full incremental monitor, the validation reference, SQLite and PyContract. They contain 75 events per implementation and are analyzed separately from the original benchmark.

OpenAI ChatGPT and Codex assisted with source recovery, implementation, reference construction, experiment orchestration, and analysis. Deterministic programs produced every scored decision, and the retained inputs and programs support replay of those decisions.

OpenAI ChatGPT and Codex also assisted with figure preparation and verification code. Figure~\ref{fig:mismatches} plots complete-decision and active-episode mismatches from saved outputs for the 987 held-out base traces. The accompanying script recomputes these counts, checks the plotted SVG labels and coordinates, and exports the retained SVGs with Inkscape 1.4. The original SVG metadata identifies Matplotlib 3.10.8.

\subsection{Comparisons and outcomes}
The original offline evaluation ran under Python 3.12.14 on Linux x86-64. Each method received the same ordered packets, time-unit declarations, base severity and policy. The base-only analysis gives each of the 987 designs equal weight. To reproduce the primary endpoint, let $a_{ij}$ and $g_{ij}$ denote the output and reference decision tuples at event $j$ of trace $i$. If $n_i$ is the trace length and $N$ the number of traces, the number of exact traces is $\sum_{i=1}^{N}\prod_{j=1}^{n_i}\mathbf{1}[a_{ij}=g_{ij}]$. A single disagreement thus changes a trace's score even if all other events agree.
The full monitor and independent retrospective reference implement the same contract. Four reduced implementations isolate different requirements: timer-only processing removes version resets and observation/suppression binding; no-evidence-binding preserves resets and timing but accepts obsolete observation versions; no-suppression-binding accepts suppressions from other versions; no-freshness removes the stale-age check while retaining rejection of future-dated observations. Every method receives identical events, policy parameters, and recovered severity. A snapshot comparator directly reports the inherited Boolean finding at each event, separately characterizing the difference between snapshot and lifecycle interfaces.

Primary agreement requires equality of every $(\text{state},\text{active},\text{severity})$ tuple in the complete trace. Event agreement uses that same tuple. Missed active decisions have reference activity true and output activity false; unwarranted active decisions reverse those values. Episode agreement compares the entire list of half-open active event-index intervals, ending at trace length for an open final episode. Severity and policy-reason agreement are separately counted. Execution failures and unresolved states are retained in the record. The fixed corpus is fully enumerated, so the endpoints are exact descriptive counts with explicit denominators. The primary analysis counts each base design once.

\subsection{Additional implementation and fixture checks}
These checks followed the primary evaluation and use separately frozen protocols. The SQL and PyContract implementations receive the original corpus, policy parameters and categorical family severities. They test execution of the same contract through different computational representations.

\paragraph{Recursive SQL implementation}
The SQL implementation expresses the complete policy as a recursive query in SQLite 3.53.1. Python validates and loads the inputs and compares the outputs; SQLite computes the transitions. The query, driver, reference hashes and analysis addendum were frozen before this implementation was scored.

\paragraph{Established monitoring engine}
A retrospective comparison encoded the complete experimental policy in PyContract, a runtime-verification library introduced by Dams et al. for event streams and parameterized state machines \cite{dams2022pycontract,pyrv2025pycontract}. The executed upstream snapshot was pinned to commit \texttt{e378eed38699bcff33e2ae71426e6b161690cdfa}; its source files were verified against Git blob hashes. The policy specification used four data states: Dormant for an inactive episode without a qualification run, Qualifying for an inactive episode with a positive run, Raised for an active episode outside clearing, and Recovering for an active episode in a clearing run. Guards check evidence quality, version changes, freshness and suppression. The unmodified PyContract engine dispatched each event and maintained the active state vector. The specification imported neither the proposed monitor nor the reference oracle.

The baseline protocol, encoding, upstream source, inputs and references were frozen before scored execution. Per-event logging records calls to the upstream dispatcher. An integration check temporarily disables each of three upstream dispatch methods, then restores it, to verify that the imported engine participates in producing the output.

\paragraph{Externally specified suppression outcomes}
A supplementary conformance check used expected suppression outcomes from the versioned Alertmanager test suite. The first four \code{Mutes} assertions in \code{TestSilencer}, version 0.28.1, specify the outcomes for an absent silence, a nonmatching exact matcher, a matching exact matcher, and an expired silence \cite{prometheusSilencer0281}. We froze their source lines, expected Boolean outcomes, and mapping before replay. The retained source has verified Git blob \code{a65b1f6337723d3ba76140081b757088213f7d5a}; the supplement includes its Apache 2.0 license and content hashes.

The adapter mapped exact equality of the source test's \code{foo} label to equality of the review-policy context, with a fixed matching policy version. Two positive observations established an active episode before each query, and the source test's relative expiration time was preserved. Where the source retained both a matching and a nonmatching silence, the fixed-label projection retained the matching candidate. The mapped output was whether the final displayed state was \code{SUPPRESSED}. This extension was frozen separately after the primary experiment.
The study implementations execute the mapped inputs, and the scored endpoint is the externally specified suppression Boolean.

\section{Results}
\subsection{Version information and output requirements}
For each of the six projections, deleting one field without supplying its comparison prevented exact agreement on at least one member of every tested pair. All 30 frozen pairs had equal projected inputs and different required decisions. For each of the six projections, ten traces formed five heterogeneous classes, giving $B=5$ and an agreement upper bound of $5/10$ on that projection's assay. A supplementary generic class-count audit reproduced these bounds from canonical projected inputs and verified that the incremental monitor agreed with all 60 witness reference sequences. Each projection is a separate interface experiment.

The results distinguish two kinds of information loss. Observation-version loss permits an obsolete positive finding to influence a new episode. Suppression-version loss permits an obsolete deferral to hide a current warning. The second effect is invisible to a metric that compares only latent activity. Current context and current policy identities also have witnesses: their omission makes it impossible to distinguish an unchanged episode from an episode that must be reset under the specified policy.

\paragraph{Effect of merging the reason categories}
For a retrospective implementation check of the six separate comparisons, we recoded each version as an alternating token on change and each nonmatching packet field as a reserved token. Across all 4,935 frozen traces, replay reproduced every reference sequence, including its reasons. This replay uses the original monitor and examines a representation property separately from the held-out implementation comparison.
In a retrospective replay we encoded the three joint comparisons, representing all changes as context changes and all mismatches as context mismatches. All 26,430 decision tuples in 4,935 traces matched the frozen reference. Of the 585 reason differences, 510 were policy mismatches recorded as context mismatches and 75 were policy changes recorded as context changes. These are counts for this explicit decoding. The paired example above establishes why a different decoder still cannot recover both reasons in general.

\subsection{Complete decisions and active episodes}
\label{sec:agreements}
The complete monitor matched the reference on all 2,961 held-out traces and all 15,858 event tuples. Table~\ref{tab:results} separates complete-decision agreement from active-episode agreement and active-event errors for the 987 base designs. Metadata and time-unit variants are excluded from this table so each base design contributes once. All six compared implementations completed every trace without execution failures. The reference contains 1,803 unresolved decisions among 5,286 base events; these represent specified missing, conflicting, or invalid evidence states.

\begin{table}[H]
\centering\small
\caption{Held-out base designs: 987 traces and 5,286 events. Trace agreement compares state, activity, and severity. Missed/unwarranted counts concern active-event decisions. Reference-active denominator: 1,293; reference-inactive denominator: 3,993.}\label{tab:results}
\begin{tabular}{@{}lrrrr@{}}\toprule
Method & Trace agreement & Episode agreement & Missed & Unwarranted\\\midrule
Version-bound &987&987&0&0\\
Timer-only &897&927&0&192\\
No evidence binding &957&969&0&84\\
No suppression binding &957&987&0&0\\
No freshness &969&975&6&18\\
Snapshot comparator &246&246&408&1,179\\\bottomrule
\end{tabular}
\end{table}

Timer-only processing disagreed on 90 of 987 base traces and produced 192 unwarranted active decisions. Across all held-out variants these become 270 trace mismatches and 576 unwarranted active decisions. Observation-binding removal caused 30 base-trace mismatches and 84 unwarranted active events. Removing the stale-age check caused 18 base-trace mismatches, six missed active events, and 18 unwarranted active events. This ablation therefore produced errors in both directions of active-event classification.

Suppression-binding removal caused 30 base-trace mismatches and 48 event-tuple mismatches, while all 987 active-episode interval sequences remained correct. Every discrepancy in this ablation concerned displayed state and severity; the active-episode endpoint registered none of them. Figure~\ref{fig:mismatches} compares the mismatched base-trace counts for the complete and reduced lifecycle implementations.

\begin{figure}[H]
\centering\includegraphics[width=\linewidth]{trace_mismatches.pdf}
\caption{Complete-decision and active-interval mismatches across the same 987 held-out base traces. Circles count traces with a differing sequence of displayed state, activity or severity. Squares count traces with differing lists of active intervals. Removing suppression binding changes complete decision sequences on 30 traces, with all active intervals preserved.}\label{fig:mismatches}
\end{figure}

Table~\ref{tab:attributes} counts discrepancies on the same 5,286 base events. Severity describes the warning shown at an event. The reason code identifies the policy condition that selected it.
\begin{table}[H]
\centering\small
\caption{Attribute discrepancies across 5,286 events in the 987 held-out base traces. The decision tuple contains state, activity and severity. The snapshot comparator uses a single snapshot-specific reason code.}\label{tab:attributes}
\begin{tabular}{@{}lrrr@{}}\toprule
Method & Decision tuple & Severity label & Reason code\\\midrule
Version-bound & 0 & 0 & 0\\
Timer-only & 342 & 342 & 282\\
No evidence binding & 204 & 204 & 204\\
No suppression binding & 48 & 48 & 48\\
No freshness & 51 & 51 & 18\\
Snapshot comparator & 1,851 & 1,761 & 5,286\\
\bottomrule
\end{tabular}
\end{table}
The full monitor's reason and severity fields also matched the reference on all 15,858 held-out events. Reason agreement identifies the same contract condition as the reference at each event. The snapshot comparator agreed on 246 base traces, reflecting the differences between immediate Boolean reporting and the study's qualification, clearing, suppression, and revision policy. All detailed results, including mechanism and family strata, are derived from the saved per-event outputs.

\subsection{Invariance under metadata and unit changes}
Each method produced identical complete outputs for all 3,290 base-to-transformed comparisons across the full benchmark: 1,645 metadata comparisons and 1,645 time-unit comparisons. This invariance held for the complete monitor, all four reduced implementations, and the snapshot comparator. The transformation checks therefore passed even for methods with mismatched reference decisions.

\subsection{Reference and concurrent-condition checks}
Replaying eleven archived source fixtures reproduced the expected rule identifiers in all eleven cases. The taxonomy audit matched the policy-input fields of all 825 intervention representations against the frozen records after time-unit normalization. These checks connect the source adapters and described intervention grid to the inputs used in the temporal evaluation.

The incremental monitor and original reference agreed on all 1,974 development traces. The validation reference rejected all 28 malformed-input probes and preserved every original benchmark and witness output. The post-evaluation corrected-reference regressions comprised 21 valid traces and eight malformed inputs, and all 29 passed.

The full incremental monitor, the validation reference, SQLite and PyContract each matched every expected final decision and reason in the 24 concurrent-condition cases: 96/96 implementation--case checks passed. Their full output sequences also agreed across all 75 events per implementation. The checks include coincident version resets and invalid observations, clearing and escalation boundaries, and suppression at exact expiry.

\subsection{Execution in SQLite and PyContract}
SQLite reproduced all 26,430 event decisions, severities and reasons on the same 4,935 traces, with zero execution failures. It matched all 60 witness sequences and passed 21 hand-trace and 28 malformed-input checks.

The PyContract implementation received all 4,935 original traces, comprising 26,430 events, with identical policies and categorical family severities. It matched the frozen reference on \textbf{4,935/4,935 complete traces} and \textbf{26,430/26,430 event tuples}; reason, severity and active-episode intervals also agreed exactly. It recorded zero missed active decisions among 6,465 reference-active events, zero unwarranted active decisions among 19,965 reference-inactive events, and zero execution failures. The original development and held-out subsets, retrospectively replayed, yielded 1,974/1,974 and 2,961/2,961 matching traces, respectively. Separate checks passed on 60/60 witness sides, 21/21 hand-calculated valid traces and 28/28 malformed-input cases.

Per-trace records confirm 26,430 calls through the upstream PyContract event dispatcher and the corresponding data-state transitions. Disabling each of the three checked dispatch methods caused the integration check to fail; restoring the methods restored the expected output.

\subsection{Projected suppression fixtures}
All four mapped outputs reproduced the externally specified suppression expectations: false for an absent silence, false for a nonmatching exact matcher, true for a matching exact matcher, and false for an expired silence. Complete output sequences also agreed with the reference, all four episode-activation checks passed, and no execution failed.

\section{Discussion}
The versions named by a finding and by a deferral answer different questions. A finding that matches the current evidence and policy can contribute to qualification or clearing. A deferral that matches those versions can change what the reviewer sees once an episode is active. In the running audit case, accepting the deferral for \code{c0} hides the warning for \code{c1} at qualification and removes its displayed severity. Both monitors show the warning at expiry. An evaluation based on the episode interval loses the difference at qualification. The retained case is \path{audit/suppression_old_context/2/base}.

Evidence age acts earlier in the decision. A stale positive observation can sustain a qualification run, whereas a stale negative observation can sustain a clearing run. Accepting the former can raise an unwarranted warning, and accepting the latter can clear an episode too early. These transitions explain why removing the stale-age check produced active-event errors in both directions. The metadata and time-unit checks address a separate question. Even implementations that accepted obsolete or stale inputs preserved both transformations, so an evaluation of semantic agreement also needs cases that exercise the disputed condition.

The choice of interface follows from the monitor's outputs. A service with access to the current and supplied version names can send their comparisons with each event. A monitor producing decisions can combine evidence and policy into a single pair. A monitor reporting the specified reasons can use the six separate comparisons to retain the source of a reset or observation mismatch. The deferral uses a joint match and expiry because its reason identifies a suppression that actually affected the display. This makes concrete the dependence of a monitoring abstraction on the property being checked \cite{henzinger2022abstract,jevitha2024abstract}.

The same policy ran in the incremental implementation, SQLite and an established runtime-monitoring engine. Agreement across those executions locates the design choice at the boundary between the producer of evidence and the monitor. Retaining version relationships upstream supplies the distinctions the policy uses, and scoring the complete displayed decision downstream exposes an obsolete authorization even when its underlying episode is correct.

\section{Conclusion}
A version-sensitive review monitor can decide what to display from three comparisons: whether the current version pair changed, whether a supplied observation matches that pair, and whether a supplied deferral matches it. Keeping the evidence and policy comparisons separate also allows the monitor to identify the source of a reset or observation mismatch. The 30 held-out base sequences on which an obsolete deferral changed the display despite an identical episode interval show why the acceptance test must compare the warning as well as its activity. A next evaluation can fix an organization's review and deferral rules, replay its historical events in order, and compare predicted intervals, displayed warnings and reason labels with the corresponding recorded decisions.

\section{Limitations}
The lifecycle policy and temporal traces are constructed research materials. The reset, timing and suppression rules require validation against an externally established review policy and its recorded decisions. All five source-rule families share the same 329 temporal templates, so the held-out partition measures conformance across family instances. Generalization to unseen temporal patterns and deployed systems remains unmeasured. The source predicates are lexical review signals; their finding accuracy and the interpretation of their rule names are upstream of the evaluated temporal decision. Relevant-evidence equivalence is also supplied as an input and requires an upstream definition of which source changes preserve it.

The reference was authored by a separate AI agent and shares the written specification with the tested monitor. The retained hand-derived checks and independent implementations verify agreement with that contract. The SQLite and PyContract executions reuse the original traces after the primary evaluation. Their results measure behavioral conformance. Comparative runtime performance remains unmeasured. The finite-benchmark bound applies to each disclosed projected-trace distribution; causal attainability of the aggregate bound remains unestablished.

The external suppression check covers four of the eight assertions in the pinned upstream test through a fixed-label projection. Future silence activation, overlapping silences and regular-expression matching lie outside that projection. The executed programs were the study implementations; the upstream Alertmanager program was not run. Operational benefit, human review effectiveness and deployment reliability require evaluation with organizational review histories and recorded actions.

\section*{Research materials}
The accompanying reproduction archive contains the frozen temporal packets and reference sequences, monitor implementations, pinned third-party source and licenses, raw outputs, and analysis scripts. From its extracted \code{Paper7\_Reproduction} directory, run \code{python -B run\_all.py ../paper7-run}. The command verifies the preserved package, replays the frozen study, independently rescores the raw outputs, and runs the intervention-taxonomy and concurrent-condition checks. The added review checks have their own manifest. The source-provenance inventory identifies the original adapters and their archived versions. Editable figures and the literature-discovery records accompany the manuscript.

\section*{Declaration of generative AI and AI-assisted technologies in the manuscript preparation process}
During preparation of this work, the author used OpenAI ChatGPT and Codex for manuscript drafting, editorial review, and figure preparation and verification code. The author takes responsibility for the content of the article.

\bibliographystyle{unsrtnat}
\bibliography{Paper7_References}
\end{document}
